%------------------------------------------------------------------------------
% 4.5 تصميم الأمان والصلاحيات
% المالك: BE (مهمة جانت ٢.٢.٦ — 09/11 إلى 12/11/2026)
% الحجم المستهدف: 2-3 صفحات
% الغرض: ترجمة سردية الخصوصية في الأطروحة إلى آليات تصميمية قابلة للفحص —
%   اتساقاً مع خطة الاستجابة للمخاطر R-09 بصيغتها المعدلة.
%
% عناصر مطلوبة:
%   [ ] 4.5.1 مصفوفة التحكم المبني على الأدوار RBAC (tab:rbac-matrix):
%       الفاعلون الأربعة (مسؤول المنصة، مسؤول تقنية المعلومات، المدير،
%       الموظف) × حالات الاستخدام السبع — بمستويات (تنفيذ / اطلاع / محجوب)،
%       مع الإشارة إلى علاقة التعميم بين الأدوار الثلاثة المؤسسية
%   [ ] 4.5.2 إدارة الجلسات والرموز: JWT قصيرة الصلاحية، تدفق المصادقة
%       (الدخول / التجديد / الإبطال)، وتحديد معدل المحاولات (Rate Limiting)
%   [ ] 4.5.3 سجل التدقيق وتتبع الوكلاء: ماذا يُسجل لكل طلب (من / متى /
%       ماذا / أي وكيل تدخل / أي عرض ذكي استُخدم / كم استغرق) وكيفية
%       استعلامه وتصفية نتائجه (حالة الاستخدام UC-06)
%
% تحذيرات:
%   - ممنوع إدخال المصادقة الثنائية MFA/2FA في التصميم — الأطروحة تنص على
%     عدم دعمها في هذه النسخة (اتساقاً مع إصلاح تناقض R-09 في الفصل الثالث)
%   - اربط كل آلية بمخاطرها في سجل المخاطر (R-xx) لتظهر السلسلة كاملة:
%     خاطر ← استجابة ← تصميم
%------------------------------------------------------------------------------
\section{تصميم الأمان والصلاحيات}
\label{sec:security-design}
~

\subsection{مصفوفة التحكم المبني على الأدوار (\EN{RBAC})}
\label{subsec:rbac-matrix}
~

\subsection{إدارة الجلسات والرموز}
\label{subsec:auth-session}
~

\subsection{سجل التدقيق وتتبع الوكلاء}
\label{subsec:audit-design}
~
